\begin{frame}{실습 3: 틱택토의 첫 수 (15.1과 연결)}
  15.1절에서 틱택토를 \textbf{완전탐색}(549,946개 국면)으로
  풀어 ``최선이면 무승부''라는 \textbf{정확한} 답을 구했다.
  이번엔 MCTS로 --- 빈판에서 5000회 시뮬레이션($c=1.4$, 시드
  0/1/2) 후 첫 수 후보 9개:
  \vskip0.5em
  \small
  \begin{tabular}{@{}ll@{}}
    \toprule
    위치 & 승률 / 방문 점유율(세 시드 범위) \\
    \midrule
    \textbf{가운데(4)} & \textbf{0.72$\sim$0.76} / \textbf{29$\sim$41\%} (1위) \\
    네 귀퉁이 & 0.62$\sim$0.68 / 각 9$\sim$13\% \\
    나머지(가운데 끼인 변) & 0.54$\sim$0.63 / 나머지 \\
    \bottomrule
  \end{tabular}
  \vskip0.8em
  \begin{itemize}
    \item 세 시드 모두 \textbf{가운데}가 방문 점유율과 승률
      \textbf{둘 다에서 1위}
  \end{itemize}
  \vskip0.4em
  \begin{block}{같은 눈으로 읽기}
    완전탐색의 ``무승부'' = \textbf{양쪽이 모두 최선}일 때의
    결과, MCTS의 승률 = \textbf{상대가 무작위 플레이}를 할 때의
    실전 승률. 틱택토엔 ``필승'' 수가 없으므로 MCTS의 순위는
    ``\textbf{바보 상대를 상대로} 얼마나 자주 이기고
    (상대가 실수할수록 이득이 크고), 얼마나 적게 지는가''를
    반영 --- 그 관점에서도 ``가운데 $\to$ 귀퉁이 $\to$ 변''은
    정석 순서와 일치. \textbf{두 방법의 출력을 서로 다른
    질문의 답}(정확한 가치 vs 표본 기반 추천)으로 읽는 연습 =
    계획 알고리즘을 쓸 때의 기본 태도.
  \end{block}
\end{frame}
